\documentclass[10pt,a4paper]{amsart}
\usepackage[margin=19mm]{geometry}
\usepackage{amsmath,amssymb,mathtools}
\usepackage[scr=rsfs,cal=euler]{mathalfa}
\usepackage{xfrac}
\usepackage[unicode,hidelinks,bookmarks=false]{hyperref}
\newcommand{\ie}{i.e.\ }
\newcommand{\cf}{cf.\ }
\newcommand{\resp}{respectively\ }
\newcommand{\Subsection}{\subsection}
\providecommand{\nobreakdash}{\nobreak\hbox{-}}
\setlength{\emergencystretch}{2em}
\multlinegap0pt
\usepackage[shortlabels]{enumitem}
\usepackage{amssymb}
\usepackage{mathtools}
\usepackage{tikz}
\newtheorem{theorem}{Theorem}
\title{\textbf{Verifying Good Regulator Conditions for Hypergraph Observers:} \\[0.3em] \large Natural Gradient Learning from Causal Invariance via Established Theorems}
\author{Max Zhuravlev}
\date{Source: arXiv:2603.09067v1 (CC BY 4.0)}
\begin{document}
\maketitle

\noindent\emph{Source and licence.} \textbf{Verifying Good Regulator Conditions for Hypergraph Observers:} \\[0.3em] \large Natural Gradient Learning from Causal Invariance via Established Theorems. Version arXiv:2603.09067v1 (CC BY 4.0), distributed by arXiv under the Creative Commons Attribution 4.0 International licence, http://creativecommons.org/licenses/by/4.0/, as recorded in the arXiv licence metadata for this exact version and shown on the abstract page https://arxiv.org/abs/2603.09067v1. Full licence notice: \texttt{licenses/arxiv-2603.09067-CC-BY-4.0.txt}.\par\smallskip

\noindent\textit{Editorial scope.} Counted unit: the article's theorem thm:optimal\_alpha (source0.tex lines 669--691). The article's complete original proof of it is delivered with it (source0.tex lines 693--699, one \texttt{proof} environment of 7 lines). No mathematics is written, repaired or re-derived: the statement and the proof are the article's own lines, in the article's own notation, and no retained line is substituted or re-flowed.  No further source-local context is needed: every cross-reference of every retained line resolves inside the two delivered spans.  Nothing was omitted from any retained span, so the package carries no omission record.  No retained line contains a citation, so the wrapper carries no bibliography and no bibliography entry.  No retained line contains a figure environment, an image-inclusion command or drawing code, so no image is delivered and the package declares no asset dependency.  Editorial formatting only, in addition: an AMS \texttt{amsart} preamble that restates the article's own declarations and macros used by the retained lines, namely the environments \texttt{theorem} on separate counters, together with the standard packages those lines need.  The retained environments print in this document as Theorem 1 in order of appearance, whereas the article prints the same items with the numbering of its own full document.  The complete native source of the article as distributed in the arXiv e-print for this exact version is bundled unchanged as \texttt{sources/196043/source0.tex}.
\begin{theorem}[Convergence-Time Optimal $\alpha$]
\label{thm:optimal_alpha}
Let $F$ be a positive definite Fisher matrix with eigenvalues $0 < \lambda_{\min} \leq \lambda_{\max}$ and condition number $\kappa = \lambda_{\max}/\lambda_{\min}$. Define the combined metric $g(c) = F^2 + cF$ and the convergence time functional:
\begin{equation}
T(c) = \kappa(g(c)) \cdot \mu_{\max}(g(c))
\end{equation}
where $\kappa(g) = \mu_{\max}/\mu_{\min}$ is the condition number of $g$ and $\mu_{\max}$ is the largest eigenvalue. Then:

\begin{enumerate}[noitemsep]
\item If $\kappa \leq 2$: $T(c)$ is monotonically increasing on $c \geq 0$, minimized at $c = 0$ (corresponding to $\alpha = 0$, classical regime).

\item If $\kappa > 2$: $T(c)$ has a unique interior minimum at:
\begin{equation}
c^* = \lambda_{\max} - 2\lambda_{\min}
\end{equation}
giving the optimal regime parameter:
\begin{equation}
\alpha_{\mathrm{opt}} = \frac{-\Delta + \sqrt{\Delta(\Delta + 4)}}{2}, \quad \Delta = \lambda_{\max} - 2\lambda_{\min}
\end{equation}

\item At the optimum, $\kappa(g(c^*)) = 2$ and the convergence time is reduced by approximately $\kappa/4$ relative to $c = 0$.
\end{enumerate}
\end{theorem}

\begin{proof}[Proof Sketch]
The convergence time $T(c)$ combines conditioning (ratio of extreme eigenvalues) and scale (maximum eigenvalue). For the metric $g(c) = F^2 + cF$, the eigenvalues are $\mu_k(c) = \lambda_k(\lambda_k + c)$. Computing the derivative:
\begin{equation}
\frac{dT}{dc} \propto (2\lambda_{\min} - \lambda_{\max} + c)
\end{equation}
This changes sign at $c^* = \lambda_{\max} - 2\lambda_{\min}$. Second derivative analysis confirms this is a minimum. An interior optimum exists (i.e., $c^* > 0$) if and only if $\lambda_{\max} > 2\lambda_{\min}$, equivalent to $\kappa > 2$. The condition number at optimum satisfies $\kappa(g(c^*)) = 2$ by construction.
\end{proof}

\end{document}
